

\section*{\SGAInlineRefTarget{sga3:iii:section:1}1. Extensions and Cohomology}
\addcontentsline{toc}{section}{1. Extensions and Cohomology}

\subsection*{\SGAInlineRefTarget{sga3:iii:subsection:1-1}1.1.}
Let \(\mathscr C\) be a category stable under fiber
products.
Let \(S\) be an object of \(\mathscr C\), let \(G\) be a representable
\(S\)-group, and let \(F\) be an \(S\)-functor in commutative groups
on which \(G\) acts. The cohomology groups \(H^n(G,F)\) were defined
in \SGARefLink{sga3:i:subsection:5-1}{I~5.1}. Recall that they are the cohomology groups of a complex
denoted \(C^\ast(G,F)\). Writing
\[
 (G/S)^n=G\times_S\cdots\times_S G
 \qquad(n\text{ factors}),
\]
one has
\[
 C^n(G,F)=\operatorname{Hom}_S((G/S)^n,F).
\]
Since \(G\), and hence each \((G/S)^n\), is representable, one also has
\[
 C^n(G,F)=F((G/S)^n).
\]
It follows from this and from the definition of the coboundary operator
that the complex \(C^\ast(G,F)\) depends only on the restriction of
\(F\) to the full subcategory of \(\mathscr C/S\) whose objects are
the Cartesian powers of \(G\) over \(S\). Consequently:

\medskip
\noindent\SGAInlineRefTarget{sga3:iii:lemma:1-1-1}\textbf{Lemma 1.1.1.}
\textit{Let \(\mathscr C\) be a category stable under fiber products,
let \(S\) be an object of \(\mathscr C\), and let \(G\) be a
representable \(S\)-group. Denote by \(\mathscr C(G)\) the full
subcategory of \(\mathscr C/S\) whose objects are the Cartesian powers
of \(G\) over \(S\). Let \(F\) and \(F'\) be two \(S\)-functors in
commutative groups on which \(G\) acts. If \(F\) and \(F'\) have the
same restriction to \(\mathscr C(G)\), there is a canonical
isomorphism}
\[
 H^\ast(G,F)\xrightarrow{\sim}H^\ast(G,F').
\]

\subsection*{\SGAInlineRefTarget{sga3:iii:subsection:1-1-2}1.1.2. Cohomology and restriction of scalars.\protect\footnotemark}
\footnotetext{Editor's note: the heading of this paragraph has been added.}
We record another comparison result. Let \(T\to S\) be a morphism in
\(\mathscr C\). If \(F\) is a \(T\)-functor in commutative groups,
then the functor obtained by restriction of scalars (cf. \SGARefLink{sga3:ii:section:1}{Expos\'e~II,
1}),
\[
 F_1=\prod_{T/S}F,
\]
is an \(S\)-functor in commutative groups, and there is a morphism of
\(S\)-functors in groups\footnotemark
\[
 u:\prod_{T/S}
 \underline{\operatorname{Aut}}_{T\text{-gr.}}(F)
 \longrightarrow
 \underline{\operatorname{Aut}}_{S\text{-gr.}}(F_1).
\]
\footnotetext{Editor's note: explicitly, let \(S'\to S\) and
\(\alpha\in
\underline{\operatorname{Aut}}_{(T\times_S S')\text{-gr.}}
(F_{T\times_S S'})\). For every \(U\to T\times_S S'\), let
\(\alpha_U\in
\operatorname{Aut}_{\mathrm{gr.}}(F_{T\times_S S'}(U))\) be the
element defined by \(\alpha\). Then \(u(S')(\alpha)\) is the element
\(\beta\) of
\[
 \underline{\operatorname{Aut}}_{S\text{-gr.}}(F_1)(S')
 =
 \underline{\operatorname{Aut}}_{S'\text{-gr.}}
 \left(\prod_{T\times_S S'/S'}F_{T\times_S S'}\right)
\]
such that \(\beta_{S''}=\alpha_{T\times_S S''}\) for every
\(S''\to S'\).}

Now let \(G\) be an \(S\)-functor in groups, and suppose an action of
\(G_T\) on \(F\) is given by
\[
 G_T\longrightarrow
 \underline{\operatorname{Aut}}_{T\text{-gr.}}(F).
\]
By the definition of \(\prod_{T/S}\), this gives a morphism of
\(S\)-functors in groups
\[
 G\longrightarrow
 \prod_{T/S}\underline{\operatorname{Aut}}_{T\text{-gr.}}(F).
\]
Composing with \(u\) therefore gives an action of \(G\) on
\(F_1=\prod_{T/S}F\).\footnote{Editor's note: explicitly, for every
\(S'\to S\), the action of \(G(S')\) on
\(F_1(S')=F(T\times_S S')\) is given by the action of
\(G_T(T\times_S S')=G(T\times_S S')\) on \(F(T\times_S S')\),
together with the group morphism
\(G(S')\to G(T\times_S S')\) corresponding to the projection
\(T\times_S S'\to S'\). Equivalently, the action of \(G_T\) on
\(F\) is a morphism \(G_T\times_T F\to F\). Applying
\(\prod_{T/S}\), and writing \(G_1=\prod_{T/S}G_T\), gives, by
\SGARefLink{sga3:ii:remark:1-2}{II~1.2}, a morphism \(G_1\times_S F_1\to F_1\), hence an action of
\(G_1\) on \(F_1\). The action of \(G\) is obtained through the
natural morphism \(G\to G_1\), which corresponds by adjunction to
\(\operatorname{id}_G\).}

\medskip
\noindent\SGAInlineRefTarget{sga3:iii:lemma:1-1-2}\textbf{Lemma 1.1.2.}
\textit{Under the preceding assumptions, there is a canonical
isomorphism}
\[
 H^\ast\!\left(G,\prod_{T/S}F\right)
 \simeq H^\ast(G_T,F).
\]
Indeed, by the definition of cohomology, the standard complexes are
canonically isomorphic.

\subsection*{\SGAInlineRefTarget{sga3:iii:subsection:1-2}1.2. Lifting group morphisms.\protect\footnotemark}
\footnotetext{Editor's note: this heading has been added.}
Following the general principles, we make the following definition.

\medskip
\noindent\SGAInlineRefTarget{sga3:iii:definition:1-2-1}\textbf{Definition 1.2.1.}
\textit{Let}
\[
 1\longrightarrow M\xrightarrow{u}E\xrightarrow{v}G
\]
\textit{be a sequence of morphisms of \(\widehat{\mathscr C}\)-groups.
It is called \emph{exact} if the following equivalent conditions hold:}
\begin{enumerate}
 \item[(i)] \textit{for every \(S\in\operatorname{Ob}\mathscr C\),
 the following sequence of ordinary groups is exact:}
 \[
  1\longrightarrow M(S)\xrightarrow{u(S)}E(S)
   \xrightarrow{v(S)}G(S);
 \]

 \item[\SGAInlineRefTarget{sga3:iii:definition:1-2-1:item:ii}(ii)] \textit{for every object \(H\) of
 \(\widehat{\mathscr C}\), the following sequence of ordinary groups
 is exact:}
 \[
  1\longrightarrow\operatorname{Hom}(H,M)
   \xrightarrow{u(H)}\operatorname{Hom}(H,E)
   \xrightarrow{v(H)}\operatorname{Hom}(H,G).
 \]
\end{enumerate}

Taking \(H=G\) in \SGARefLink{sga3:iii:definition:1-2-1:item:ii}{(ii)}, one sees that the set of sections of \(v\)
(which need not preserve the group structures) is either empty or
principal homogeneous under \(\operatorname{Hom}(G,M)\). Suppose that
it is nonempty, and let
\[
 s:G\longrightarrow E
\]
be a section of \(v\). For every \(S\in\operatorname{Ob}\mathscr C\)
and every \(x\in G(S)\), the element \(s(x)\in E(S)\) defines an
inner automorphism of \(E_S\) that normalizes \(M_S\) (more precisely,
the image of \(M_S\) under \(u_S\)), and hence an automorphism of
\(M_S\).

\medskip
\noindent\SGAInlineRefTarget{sga3:iii:scholium:1-2-1-1}\textbf{Scholium 1.2.1.1.}
\textit{If \(M\) is commutative, one sees set-theoretically that this
automorphism is independent of the chosen section and depends only on
\(x\), and that it depends multiplicatively on \(x\). In summary, every
exact sequence}
\[
 (E)\qquad
 1\longrightarrow M\xrightarrow{u}E\xrightarrow{v}G
\]
\textit{in which \(M\) is commutative and \(v\) admits a section
determines a morphism of \(\widehat{\mathscr C}\)-groups}
\[
 G\longrightarrow
 \underline{\operatorname{Aut}}_{\widehat{\mathscr C}\text{-gr.}}(M),
\]
\textit{called the \emph{action of \(G\) on \(M\) defined by the
extension \((E)\)}.}

\medskip
\noindent\SGAInlineRefTarget{sga3:iii:definition:1-2-1-2}\textbf{Definition 1.2.1.2.}
We saw in \SGARefLink{sga3:i:proposition:2-3-7}{I~2.3.7} that \(v\) admits a section that is a morphism of
\(\widehat{\mathscr C}\)-groups if and only if the extension \((E)\)
is isomorphic, as an extension, to the semidirect product of \(M\) by
\(G\) for the preceding action. Such a section of \(v\) is called a
\emph{section of the extension \((E)\)}, or simply a \emph{section of
\((E)\)}.

If \(s\) is a section of \((E)\) and
\[
 m\in\Gamma(M)\simeq
 \operatorname{Ker}\bigl(\Gamma(E)\to\Gamma(G)\bigr)
\]
(for the definition of \(\Gamma\), see \SGARefLink{sga3:i:subsection:1-2}{I~1.2}), then the morphism
\(G\to E\) defined by
\[
 x\longmapsto u(m)s(x)u(m)^{-1}
\]
is also a section of \((E)\). It is called
the section deduced from \(s\) by the inner automorphism defined by
\(m\), or by \(u(m)\).

\medskip
\noindent\SGAInlineRefTarget{sga3:iii:lemma:1-2-2}\textbf{Lemma 1.2.2.}
\textit{Let}
\[
 (E)\qquad
 1\longrightarrow M\xrightarrow{u}E\xrightarrow{v}G
\]
\textit{be an exact sequence of \(\widehat{\mathscr C}\)-groups in
which \(M\) is commutative and \(v\) admits a section. Let \(G\) act
on \(M\) in the manner defined by \((E)\).}
\begin{enumerate}
 \item[\SGAInlineRefTarget{sga3:iii:lemma:1-2-2:item:i}(i)] \textit{The extension \((E)\) canonically defines a class
 \(c(E)\in H^2(G,M)\), whose vanishing is necessary and sufficient for
 the existence of a section of \((E)\).}

 \item[\SGAInlineRefTarget{sga3:iii:lemma:1-2-2:item:ii}(ii)] \textit{If \(c(E)=0\), the set of sections of \((E)\) is
 principal homogeneous under \(Z^1(G,M)\), and the set of sections of
 \((E)\), modulo the action of the inner automorphisms defined by the
 elements of \(\Gamma(M)\), is principal homogeneous under
 \(H^1(G,M)\).}

 \item[\SGAInlineRefTarget{sga3:iii:lemma:1-2-2:item:iii}(iii)]\footnote{Editor's note: this item has been added as the
 counterpart of \SGARefLink{sga3:iii:proposition:1-2-4:item:iii}{1.2.4(iii)}. The assertion is valid for every section
 of \(v\); see the proof.}
 \textit{Let \(s\) be a section of \((E)\). The set of conjugates of
 \(s\) by the inner automorphisms defined by \(\Gamma(M)\) is in
 bijection with \(\Gamma(M)/H^0(G,M)\).}
\end{enumerate}

The proof is exactly the same as for ordinary groups; starting from a
section of \(v\) ensures the functoriality of the set-theoretic
computations. We briefly indicate the principal steps.

\smallskip
\noindent\textbf{(a)}
To every section \(s\) of \(v\), associate the morphism
\[
 D_s:G\times G\longrightarrow M
\]
defined set-theoretically by
\[
 u(D_s(x,y))=s(xy)s(y)^{-1}s(x)^{-1}.
\]
The following computation shows that \(D_s\) is a
2-cocycle. By the definition of the differential
in the standard complex (\SGARefLink{sga3:i:subsection:5-1}{I~5.1}),
\[
\begin{aligned}
 (\partial^2D_s)(x,y,z)
  ={}&\bigl(s(x)D_s(y,z)s(x)^{-1}\bigr)
     \cdot D_s(x,y)^{-1}\cdot D_s(xy,z)^{-1}\cdot D_s(x,yz).
\end{aligned}
\]
Substituting the definition of \(D_s\) into this formula gives,
without using commutativity,
\[
 D_s\in Z^2(G,M).
\]

\smallskip
\noindent\textbf{(b)}
If \(s\) and \(s'\) are two sections of \(v\), there is a morphism
\(f:G\to M\) such that \(s(x)=f(x)s'(x)\). Then
\[
\begin{aligned}
 D_{s'}(x,y)
 &=f^{-1}(xy)D_s(x,y)s(x)f(y)s(x)^{-1}f(x)\\
 &=D_s(x,y)\cdot f^{-1}(xy)
   \cdot\bigl(s(x)f(y)s(x)^{-1}\bigr)\cdot f(x),
\end{aligned}
\]
and therefore
\[
 D_{s'}=D_s\cdot\partial^1f.
\]
This shows that the class of \(D_s\) in \(H^2(G,M)\) is independent
of the chosen section \(s\) of \(v\); it is the class \(c(E)\) of the
extension \((E)\).\footnote{Editor's note: the preceding sentence has
been added.}

\smallskip
\noindent\textbf{\SGAInlineRefTarget{sga3:iii:lemma:1-2-2:item:c}(c)}
Let \(s\) and \(s'\) be two sections of \(v\), and let
\(m\in\Gamma(M)\). The equality
\[
 s(x)=m^{-1}s'(x)m
 \qquad(x\in G(S),\ S\in\operatorname{Ob}\mathscr C)
\]
is equivalent to
\[
 s(x)=m^{-1}s'(x)m\,s'(x)^{-1}s'(x),
 \qquad\text{that is,}\qquad
 s=\partial^0m\cdot s'.
\]
In particular, the stabilizer of \(s\) in \(\Gamma(M)\) is the
subgroup of elements \(m\) such that
\(\partial^0m=e_M\), namely \(H^0(G,M)\). This already proves \SGARefLink{sga3:iii:lemma:1-2-2:item:iii}{(iii)}.

\smallskip
\noindent\textbf{(d)}
The argument is now familiar.\footnote{Editor's note: what follows has
been expanded.} Let \(s_0\) be an arbitrary section of \(v\). There
exists a section \(s\), necessarily of the form \(s=f\cdot s_0\), that
is a group morphism, that is, satisfies \(D_s=0\), if and only if
\[
 (D_{s_0})^{-1}=\partial^1f,
\]
for some \(f:G\to M\), or equivalently if and only if \(c(E)=0\). This
proves \SGARefLink{sga3:iii:lemma:1-2-2:item:i}{(i)}.

In that case, the sections of \((E)\) are precisely the sections
\(s'=h\cdot s\), where \(h:G\to M\) satisfies
\(\partial^1h=0\), that is, \(h\in Z^1(G,M)\). Moreover, by \SGARefLink{sga3:iii:lemma:1-2-2:item:c}{(c)}, two
such sections \(h_1\cdot s\) and \(h_2\cdot s\) are conjugate under
\(\Gamma(M)\) if and only if \(h_1\) and \(h_2\) have the same image
in \(H^1(G,M)\). This proves \SGARefLink{sga3:iii:lemma:1-2-2:item:ii}{(ii)}.

Still let
\[
 (E)\qquad
 1\longrightarrow M\xrightarrow{u}E\xrightarrow{v}G
\]
be an exact sequence of \(\widehat{\mathscr C}\)-groups with \(M\)
commutative, and let
\[
 f:H\longrightarrow G
\]
be a morphism of \(\widehat{\mathscr C}\)-groups.

Consider \(E_f=H\times_G E\). It is a
\(\widehat{\mathscr C}\)-group, and the projection
\(v_f:E_f\to H\) is a morphism of
\(\widehat{\mathscr C}\)-groups; the same holds for
\(e_f:E_f\to E\). Sending \(M\) into \(E\) by \(u\) and into \(H\)
by the unit morphism also defines a morphism of
\(\widehat{\mathscr C}\)-groups \(u_f:M\to E_f\). We have therefore
constructed the following commutative diagram of
\(\widehat{\mathscr C}\)-groups.

\SGAThreeNativeDiagram[0.82\textwidth]
  {native_diagrams/exp3/diagram_028_exp3_p027_pullback_extension_diagram.tex}
  {The pullback extension \((E_f)\) and its morphism to \((E)\).}
  {fig:exp3-pullback-extension}
  {LNM 151, pp.~105--107}

It follows immediately that:

\medskip
\noindent\SGAInlineRefTarget{sga3:iii:lemma:1-2-3}\textbf{Lemma 1.2.3.}
\begin{enumerate}
 \item[(i)] \textit{The sequence \((E_f)\) is exact.}

 \item[(ii)] \textit{The map \(s\mapsto e_f\circ s=f'\) gives a
 bijective correspondence between the morphisms}
 \[
  s:H\longrightarrow E_f
 \]
 \textit{such that \(v_f\circ s=\operatorname{id}\), that is, the
 sections of \(v_f\), and the morphisms}
 \[
  f':H\longrightarrow E
 \]
 \textit{such that \(v\circ f'=f\), that is, the morphisms \(f'\)
 lifting \(f\).}

 \item[(iii)] \textit{Under this correspondence, the sections of
 \((E_f)\) correspond to the group morphisms \(f'\) lifting \(f\).}
\end{enumerate}

Applying \SGARefLink{sga3:iii:lemma:1-2-2}{Lemma~1.2.2} to the extension \((E_f)\), and taking \SGARefLink{sga3:iii:lemma:1-2-3}{1.2.3}
into account, gives the following proposition, which formally contains
\SGARefLink{sga3:iii:lemma:1-2-2}{Lemma~1.2.2}.

\medskip
\noindent\SGAInlineRefTarget{sga3:iii:proposition:1-2-4}\textbf{Proposition 1.2.4.}
\textit{Let}
\[
 (E)\qquad 1\longrightarrow M\longrightarrow E\xrightarrow{v}G
\]
\textit{be an exact sequence of \(\widehat{\mathscr C}\)-groups with
\(M\) commutative. Let \(f:H\to G\) be a morphism of
\(\widehat{\mathscr C}\)-groups, and suppose that it lifts to a
morphism, not necessarily a group morphism, \(f':H\to E\). Let \(H\)
act on \(M\) by the composite morphism, which is multiplicative and
independent of the choice of \(f'\),}
\[
 H\xrightarrow{f'}E\xrightarrow{\operatorname{int}}
 \underline{\operatorname{Aut}}_{\widehat{\mathscr C}\text{-gr.}}(M).
\]
\begin{enumerate}
 \item[(i)] \textit{The morphism \(f\) canonically defines a class
 \(c(f)\in H^2(H,M)\), whose vanishing is necessary and sufficient for
 the existence of a morphism of \(\widehat{\mathscr C}\)-groups}
 \[
  f':H\longrightarrow E
 \]
 \textit{lifting \(f\).}

 \item[(ii)] \textit{If \(c(f)=0\), the set of morphisms of
 \(\widehat{\mathscr C}\)-groups \(f'\) lifting \(f\), modulo the
 action of the inner automorphisms defined by elements of \(\Gamma(M)\)
 (that is, by elements \(m\in\Gamma(E)\) such that \(v(m)=e\)), is
 principal homogeneous under \(H^1(H,M)\).}

 \item[\SGAInlineRefTarget{sga3:iii:proposition:1-2-4:item:iii}(iii)] \textit{If \(f':H\to E\) is a group morphism lifting
 \(f\), the set of transforms of \(f'\) by the inner automorphisms
 defined by the elements of \(\Gamma(M)\) is isomorphic to}
 \[
  \Gamma(M)/\Gamma(M^H)=\Gamma(M)/H^0(H,M).
 \]
\end{enumerate}

